\documentstyle[%


righttag%


,11pt%


,amssymb%


,russian%               remove in Eng.


,izvru%                 Set izven% in Eng.


% ,izvru-ks             for brief communications


]{amsart}





%





\newcommand{\Int}{\operatorname{int}}


\newcommand{\tts}[1]{}


\renewcommand{\baselinestretch}{0.97}





%\hoffset=-15mm%                  For Eng.


\setcounter{page}{37}


\god{2002}


\nomer{2~(477)} %                        Remove in Eng.


%\nomer{Vol.\,46, No.\,1}%               Set in Eng.


%\str{\pageref{begin}--\pageref{end}}%  Set in Eng.


\udk{517.929}





\author{\it Г.Г.\,Исламов}


% NB:   ^^ remove in Eng.


\title{О допустимых помехах


линейных управляемых систем}





\date{}


%\pagestyle{myheadings}%         Set in Eng.


\pagestyle{plain} %              Remove in Eng.


\begin{document}


%\thispagestyle{plain}%          Set in Eng.


\maketitle


\label{begin}








Пусть управляемый процесс $x:[0,T]\to R^n$ описывается линейным уравнением


\begin{equation}


\Dot x(t) + A(t)x(t) = \int_0^t K(t,s)\Dot x(s)\, ds +


B(t)u(t) + \nu(t), \quad t\in [0,T]. 


\end{equation}


Предполагается, что $n\times n$-матрица $A(t)$, $n\times k$-матрица


$B(t)$ и $n$-мерный вектор-столбец $\nu(t)$ (``помеха'') имеют


квадратично суммируемые на $[0,T]$ вещественные компоненты, ядро $K(t,s)$ ---


$n\times n$-матричная функция с квадратично суммируемыми на


$0\le s\le t\le T$ вещественными элементами, $k$-мерный вектор-столбец


$u(t)$ (``допустимое управление'') имеет измеримые на $[0,T]$ вещественные


компоненты, причем при почти всех $t\in[0,T]$ значения $u(t)$ принадлежат


компакту $\Omega$ координатного пространства $R^k$.





Под решением системы \thetag1 понимается $n$-мерный вектор-столбец


$x(t)$ с производной $\Dot x(t)$, имеющей квадратично суммируемые на


$[0,T]$ компоненты.





Помеха $\nu(t)$ называется допустимой, если найдется такое допустимое


управление $u(t) = u_\nu(t)$, при котором у системы (1) существует


решение $x(t)$, удовлетворяющее векторному неравенству


\begin{equation}


\int_0^T \Phi(s)\Dot x(s)\,ds + \Psi x(0) \ge \beta. 


\end{equation}


Здесь $\Phi(s)$ --- $m\times n$-матрица с квадратично суммируемыми на


$[0,T]$ элементами, $\Psi$ --- постоянная $m\times n$-матрица,


$\beta$ --- $m$-мерный вектор-столбец. Сравнение векторов в (2)


покомпонентное.





Будем говорить, что управление $u_\nu(t)$ компенсирует помеху $\nu(t)$.





Критерий разрешимости уравнения (1) с ограничением (2)


при отсутствии управления ($u(t)\equiv 0$) анонсирован в \cite{1}


и доказан для точечных неравенств специального вида в \cite{2}.


Подходы к исследованию разрешимости задач с ограничениями типа неравенств


для различных классов функционально-дифференциальных уравнений


см. в \cite{3}--\cite{6}.





Описание множества всех помех $\nu(t)$ системы (1), компенсируемых


заданным допустимым управлением $u(t)$, дает





\begin{thm}


Пусть $Z(s)$ есть единственное $m\times n$-матричное решение


с квадратично суммируемыми на $[0,T]$ элементами уравнения


Вольтерра


\begin{equation}


Z(s) = \int_s^T Z(\tau)[K(\tau,s) - A(\tau)]\,d\tau + \Phi(s),


\quad s\in[0,T]. 


\end{equation}


Для того чтобы система $(1)$ имела решение $x(t)$, удовлетворяющее


$(2)$, необходимо и достаточно выполнения неравенства


\begin{equation}


\lambda(\beta - \int_0^T Z(s)\{B(s)u(s) + \nu(s)\}\,ds) \le 0 


\end{equation}


для каждой $m$-мерной вектор-строки $\lambda$ с неотрицательными


компонентами, являющейся решением однородной системы линейных


алгебраических уравнений


\begin{equation}


\lambda(\Psi - \int_0^T Z(s)A(s)\,ds) = 0.


\end{equation}


\end{thm}





\begin{pf}


Обозначим $f(t) = B(t)u(t) + \nu(t)$.


Подстановка $x(t) = q + \int\limits_0^t y(s)\,ds$, $y(t) = \Dot x(t)$


в систему (1) приводит к интегральному уравнению Вольтерра


$$


y(t) = \int_0^t[K(t,s) - A(t)]y(s)\,ds + f(t) - A(t)q, \quad t\in[0,T].


$$


Единственное квадратично суммируемое решение этого уравнения можно


выразить через резольвенту $R(t,\tau)$ ядра $K(t,\tau) - A(t)$


в виде


$$


y(t) = \int_0^t R(t,\tau)[f(\tau) - A(\tau)q]\,d\tau +f(t) - A(t)q,


\quad t\in[0,T].


$$


Замена $\Dot x(s) = y(s),\ x(0) = q$ в неравенстве (2) дает


$$


\int_0^T \Phi(s)\bigg(\int_0^s R(s,\tau)[f(\tau) - A(\tau)q]\,d\tau +


f(s) - A(s)q \bigg)\,ds + \Psi q \ge\beta.


$$


После элементарных преобразований левой части этого неравенства найдем,


что постоянный $n$-мерный вектор-столбец $q$ следует выбирать из условия


справедливости векторного неравенства


\begin{equation}


\bigg[\Psi - \int_0^T Z(s)A(s)\,ds\bigg]q \ge\beta - \int_0^T Z(s)f(s)\,ds, 


\end{equation}


где


$Z(s) = \int\limits_s^T \Phi(\tau)R(\tau,s)\,d\tau + \Phi(s)$, $s\in[0,T]$.


Из последнего представления видим, что $Z(s)$ совпадает с единственным


решением интегрального уравнения (3), т.\,к. $R(\tau,s)$ есть


резольвента ядра $K(\tau,s) - A(\tau)$.





В силу обобщения теоремы Фредгольма (напр., \cite{7}, с.\,266)


векторное неравенство (6) относительно $q\in R^n$ имеет решение


в том и только том случае, когда скалярное неравенство


\begin{equation}


\lambda(\beta - \int_0^T Z(s)f(s)\,ds) \le 0 


\end{equation}


имеет место для всякой $m$-мерной вектор-строки $\lambda$ с


неотрицательными компонентами, являющейся решением однородной системы


линейных алгебраических уравнений (5).


Учитывая, что $f(s) = B(s)u(s) + \nu(s)$, видим, что (7) совпадает


с (4).


\end{pf}





\begin{note}


В действительности неравенство (4) достаточно проверить на конечном


числе образующих $\lambda^1,\dotsc,\lambda^l$ конуса неотрицательных решений


системы (5). Для отыскания этих образующих можно применить


вычислительную схему Н.В.\,Черниковой (\cite{8}, с.\,255) либо


воспользоваться модификацией симплекс-метода.


\end{note}





Следующее утверждение показывает, что допустимые помехи можно компенсировать


управлениями, удовлетворяющими некоторому принципу максимума. Из этого


принципа видно, что при формировании компенсирующего управления $u_\nu(t)$


для заданной допустимой помехи $\nu(t)$ можно ограничиться лишь крайними


точками компакта $\Omega$.





\begin{thm}


Пусть $\nu(t)$ есть допустимая помеха для задачи управления


$(1)$--$(2)$. Тогда для некоторого неотрицательного решения


$\alpha = \lambda_\nu$ системы $(5)$ найдется допустимое управление


$u^*(t),$ удовлетворяющее при почти всех $s\in[0,T]$ принципу максимума


\begin{equation}


\alpha Z(s)B(s)u^*(s) = \max_{u\in\Omega}\alpha Z(s)B(s)u 


\end{equation}


и компенсирующее помеху $\nu(t)$.


\end{thm}





\begin{pf}


Пусть $\lambda^1,\dotsc,\lambda^l$ ---


образующие конуса неотрицательных решений системы (5). В силу


теоремы 1 компенсирующее помеху $\nu(t)$ допустимое управление


$u(t) = u_\nu(t)$ удовлетворяет системе неравенств


\begin{equation}


\lambda^i(\beta - \int_0^T Z(s)[B(s)u(s) + \nu(s)]\,ds) \le 0,


\quad i=\overline{1,l}. 


\end{equation}


Рассмотрим конечномерные множества


\begin{gather*}


V = \{g\in R^m: \lambda^i g \le 0,\ i=\overline{1,l}\}, \\


W_\nu = \bigg\{g\in R^m: g = \beta - \int_0^T Z(s)\nu(s)\,ds -


\int_0^T Z(s)B(s)u(s)\,ds\bigg\},


\end{gather*}


где $u$ пробегает все измеримые управления c $u(s)\in\Omega$ при почти всех


$s\in[0,T]$.





Пересечение $V\cap W_\nu$ не пусто, т.\,к. $\nu(s)$ --- допустимая


помеха. Кроме того, $V$ и $W_\nu$ выпуклы и замкнуты, причем $V$ ---


конус с непустой внутренностью $\Int V$, а $W_\nu$ --- компакт.


Для первого множества это очевидно, а для установления свойств второго


достаточно воспользоваться (напр., \cite{9}, с.\,351) теоремой


А.А.\,Ляпунова, из которой следует выпуклость, замкнутость и ограниченность


в $R^m$ множества значений интеграла $\int\limits_0^T Z(s)B(s)u(s)\,ds$, когда


$u$ пробегает множество всех измеримых вектор-функций с $u(s)\in\Omega$


при почти всех $s\in[0,T]$, $\Omega$ --- заданный компакт в $R^k$.





Пусть направление $\alpha \ne 0$ является неотрицательной линейной


комбинацией образующих $\lambda^1,\dotsc,\lambda^l$ конуса неотрицательных


решений системы (5). Найдется граничная точка $g^*$ компакта


$V\cap W_\nu$ такая, что $\alpha g^* \le \alpha g$ при всех


$g\in V\cap W_\nu$. Допустим, что $h\in \Int V$ и 


$g^* + h\in V\cap W_\nu$. Тогда $\alpha(g^* + h) \ge \alpha g^*$


и, значит, $\alpha h \ge 0$. Однако для $h\in\Int V$ имеем


$\lambda^i h < 0$, $i=\overline{1,l}$, и, следовательно, $\alpha h < 0$.


Возникшее противоречие показывает, что $(g^* + \Int V)\cap V


\cap W_\nu = \emptyset$. Так как $g^*\in V$, то


$(g^* + \Int V)\cap V = (g^* + \Int V)$. Следовательно,


$(g^* + \Int V)\cap W_\nu = \emptyset$. Согласно первой теореме


отделимости (\cite{9}, с.\,124) открытое множество $g^* + \Int V$


и выпуклый компакт $W_\nu$ можно разделить гиперплоскостью


$\lambda_\nu g = \lambda_\nu g^*$.





Будем считать, что при всех $g\in W_\nu$


\begin{equation}


\lambda_\nu g^* \le \lambda_\nu g. 


\end{equation}


Тогда $\lambda_\nu(g - g^*)\le 0$ для всех $g\in(g^* + \Int V)$


и, значит, $\lambda_\nu h \le 0$ при всех $h\in V$. По теореме Фаркаша


(\cite{7}, с.\,247) $\lambda_\nu$ есть неотрицательная линейная комбинация


образующих $\lambda^1,\dotsc,\lambda^l$ конуса неотрицательных решений


системы (5), т.\,е. $\lambda_\nu$ есть неотрицательное решение


системы (5).





Элементу $g^*\in W_\nu$ соответствует допустимое управление $u^*(t)$


такое, что


$$


g^* = \beta - \int_0^T Z(s)\nu(s)ds - \int_0^T Z(s)B(s)u^*(s)\,ds.


$$


Для этого управления имеем $\lambda^ig^* \le 0$, $i=\overline{1,l}$,


т.\,е. выполнено (9) при $u(s) = u^*(s)$ и, значит, неравенство


(4). Утверждается, что $u^*(s)$ удовлетворяет принципу максимума


(8) при $\alpha = \lambda_\nu$. В противном случае согласно


(\cite{10}, с.\,175) выберем допустимое управление $\wt u(s)$,


удовлетворяющее принципу максимума


$$


\lambda_\nu Z(s)B(s)\wt u(s) = \max_{u\in\Omega}\lambda_\nu Z(s)B(s)u,


\quad s\in[0,T].


$$


Тогда на множестве положительной меры имеем


$$


\lambda_\nu Z(s)B(s)u^*(s) < \lambda_\nu Z(s)B(s)\wt u(s).


$$


Отсюда для элемента $\wt g = \beta - \int\limits_0^T Z(s)\nu(s)ds -


\int\limits_0^T Z(s)B(s)\wt u(s)\,ds$ из $W_\nu$ получим


$\lambda_\nu \wt g < \lambda_\nu g^*$, что противоречит (10).


\end{pf}





Из доказанных выше теорем 1 и 2 вытекает следующий критерий допустимости


помехи $\nu(t)$ для задачи управления (1)--(2).


\medskip








{\bf Следствие.}


Для того чтобы помеха $\nu(t)$ была допустимой для задачи управления


(1)--(2), необходимо и достаточно, чтобы для некоторого


неотрицательного решения $\alpha=\lambda_\nu$ системы (5) нашлось


управление $u^*(s)$, удовлетворяющее при почти всех $s\in [0,T]$


принципу максимума (8) и при $u(s) = u^*(s)$ неравенству


(4) для каждого неотрицательного решения $\lambda$ системы (5).


\medskip





\begin{note}


Пусть существует предел $\Phi(+0)$ и


$K(\tau,+0) = 0$ при почти всех $\tau\in [0,T]$. Тогда из (3)


находим, что на множестве полной меры существует предел $Z(+0)$,


равный $\Phi(+0) - \int\limits_0^T Z(\tau)A(\tau)\,d\tau$, и, значит, система


алгебраических уравнений (5) запишется в виде


$$


\lambda(\Psi - \Phi(+0) + Z(+0)) = 0. 


$$


\end{note}





\begin{note}


Если $\Phi(s)\equiv \Phi$ --- постоянная


$m\times n$-матрица, то $m\times n$-матричное решение уравнения (3)


представимо в виде $Z(s) = \Phi Y(s)$, где $Y(s)$ --- единственное


$n\times n$-матричное решение уравнения


$$


Y(s) = \int_s^T Y(\tau)[K(\tau,s) - A(\tau)]\,d\tau + E_n


$$


с единичной матрицей $E_n$ порядка $n$.


\end{note}





\begin{note}


Из представления $x(T) = \int\limits_0^T\Dot x(s)\,ds


+ x(0)$ для абсолютно непрерывной вектор-функции видно, что, взяв в


неравенстве (2)\; $\Phi(s)\equiv P$ и $\Psi = P + Q$, где


$P$, $Q$ --- постоянные $m\times n$-матрицы, получим векторное краевое


неравенство $Px(T) + Qx(0)\ge\beta$. 


\end{note}








\begin{thebibliography}{99}





\bibitem{1}


Исламов Г.Г. {\em О разрешимости интегро-дифференциальных уравнений


с краевыми неравенствами} // Тез. докл. VII Четаевской конф.


Аналитическая механика. Устойчивость и управление движением. -- Казань,


1997. -- С.\,142.


\bibitem{2}


Исламов Г.Г. {\em О достижимости полиэдров пространства состояний в заданные


моменты времени} // Вестн. Удмуртск. ун-та. -- Ижевск, 1999. -- Вып.\,8.


-- С.\,32--38.


\bibitem{3}


Исламов Г.Г. {\em Критерий разрешимости уравнений с краевыми неравенствами}


// Изв. отд. матем. и информатики УдГУ. -- Ижевск, 1994. -- Вып.\,2.


-- С.\,3--24.


\bibitem{4}


Исламов Г.Г. {\em О полиэдральной разрешимости системы линейных


дифференциальных уравнений} // Изв. вузов. Математика. -- 1999. -- \No\,3.


-- С.\,31--37.


\bibitem{5}


Исламов Г.Г. {\em О разрешимости уравнений с краевыми неравенствами}


// Краев. задачи. -- Пермь, 1981. -- С.\,88--90.


\bibitem{6}


Максимов В.П., Румянцев А.Н. {\em Краевые задачи и задачи импульсного


управления в экономической динамике. Конструктивное исследование}


// Изв. вузов. Математика. -- 1993. -- \No\,5. -- С.\,56--71.


\bibitem{7}


Беклемишев Д.B. {\em Дополнительные главы линейной алгебры}. Учеб. пособие.


-- М.: Наука, 1983. -- 336 с.


\bibitem{8}


Черников С.Н. {\em Линейные неравенства}. -- М.: Наука, 1968. -- 488 с.


\bibitem{9}


Алексеев В.М., Тихомиров В.М., Фомин С.В. {\em Оптимальное управление}.


Учеб. пособие. -- М.: Наука, 1979. -- 432 с.


\bibitem{10}


Ли Э.Б., Маркус Л. {\em Основы теории оптимального управления}.


-- М.: Наука, 1972. -- 476 с.





\end{thebibliography}





\vskip 1cm





\post{\it Удмуртский государственный}{\it Поступила}


\post{\it университет}{03.03.2000}





\label{end}


\end{document}


